\clearpage\fancyhead[L]{\small GAUSS AN BODE.}\hypertarget{native-407-part-middle}{}\label{native-407-part-middle}
\begin{center}\large GAUSS AN BODE.\end{center}
\begin{center}\hrule\vspace{4pt}\fontsize{10}{12}\selectfont Astronomisches Jahrbuch für 1821 Seite 232. Berlin 1818.\par\vspace{4pt}\hrule\vspace{4pt}\end{center}
\par\indent Beobachtete Sternbedeckungen, unterm 8. März 1818 eingesandt.
\begin{center}\fontsize{9}{11}\selectfont\setlength{\tabcolsep}{3pt}
\begin{tabular}{lll}
1818 Januar 27. & $\begin{array}{rrrl}11^{\mathrm u}&25^{\mathrm m}&27^{\mathrm s}0&\text{W. Z.}\end{array}$ & Austritt $1\;\text{\AstroSymbols ♍}$.\\[3pt]
Februar 13. & $\left.\begin{array}{rrrll}6&47&24.5&\text{—}&\text{\textsc{Gauss}}\\&&24.6&\text{—}&\text{\textsc{Harding}}\end{array}\right\}$ & Eintritt $1\;A\;\text{\AstroSymbols ♉}$.\\[3pt]
\hfill — & $\begin{array}{rrrl}8&10&26.4&\text{— H.}\end{array}$ & Austritt $1\;A\;\text{\AstroSymbols ♉}$.\\[3pt]
\hfill — & $\left.\begin{array}{rrrll}7&10&21.5&\text{—}&\text{G.}\\&&21.6&\text{—}&\text{H.}\end{array}\right\}$ & Eintritt $2\;A\;\text{\AstroSymbols ♉}$.\\[3pt]
\hfill — & $\begin{array}{rrrl}8&26&2.4&\text{— G.}\end{array}$ & Austritt $2\;A\;\text{\AstroSymbols ♉}$.\\[3pt]
\hfill — & $\begin{array}{rrrl}7&11&32.4&\text{— H.}\end{array}$ & Eintritt Stern 8. Grösse.\\[3pt]
\hfill — & $\begin{array}{rrrl}11&42&27.1&\text{— H.}\end{array}$ & Eintritt 100 \emph{Lacaille}.\\[3pt]
\hfill — & $\begin{array}{rrrl}12&30&59.9&\text{— H.}\end{array}$ & Eintritt Stern 8. Grösse.\\[3pt]
Februar 17. & $\begin{array}{rrrl}9&47&28.7&\text{— H.}\end{array}$ & Eintritt Stern 6. Grösse.\\
\end{tabular}\end{center}
\clearpage\fancyhead[L]{\small COMET.}\hypertarget{native-407-part-lower}{}\label{native-407-part-lower}
\begin{center}\large COMET.\end{center}
\begin{center}\hrule\vspace{4pt}\fontsize{10}{12}\selectfont V. \textsc{Lindenau} u. \textsc{Bohnenberger} Zeitschrift für Astronomie u. verwandte Wiss.\par Band V. Seite 276..277. März und April 1818.\par\vspace{4pt}\hrule\vspace{4pt}\end{center}
\par\indent Aus den Beobachtungen des Cometen von Dr. \textsc{Olbers} vom 3. bis 28. März
habe ich dieser Tage folgende parabolische Elemente abgeleitet, die von den von
\textsc{Olbers} selbst aus seinen Beobachtungen vom 3. und 13. März und der mir noch
unbekannten Marseiller vom 18. Januar erhaltenen, nur wenig abweichen.

\begin{center}\fontsize{10}{12}\selectfont
Durchgang durch das Perihel 1818. Februar 26,8991 Göttinger Zeit.
\par\smallskip\begin{tabular}{lrr}
Logarithm. des Abstandes im Perihel & $\cdot\;\cdot\;\cdot\;\cdot\;\cdot\;\cdot\;\cdot\;\cdot$ & $0.07937$\\
Neigung der Bahn & $\cdot\;\cdot\;\cdot\;\cdot\;\cdot\;\cdot\;\cdot\;\cdot\;\cdot\;\cdot\;\cdot\;\cdot\;\cdot\;\cdot\;\cdot$ & $90^\circ\;0'\;0''$\\
Aufsteigender Knoten & $\cdot\;\cdot\;\cdot\;\cdot\;\cdot\;\cdot\;\cdot\;\cdot\;\cdot\;\cdot\;\cdot\;\cdot\;\cdot\;\cdot$ & $70\;5\;12$\\
\end{tabular}
\par\smallskip Länge des Perihels, die Bewegung als rechtläufig betrachtet $183\;22\;58$
\end{center}
\SourceLeaf{408}\par\indent Die senkrechte Lage der Bahn gegen die Ekliptik ergab sich ungezwungen aus den
Beobachtungen, wird aber wohl nach fortgesetzten Beobachtungen einige
Modification erhalten. Dann erst wird man sehen, ob die Bewegung in Beziehung
auf die Ekliptik rechtläufig oder rückläufig ist. Ich habe zur Erleichterung
der Beobachtungen aus diesen Elementen folgende Ephemeride berechnet. Die
Zeit ist $14^{\mathrm u}20^{\mathrm m}$ in Göttingen. Die Columne für die
Lichtstärke setzt erborgtes Licht voraus, $=1$ in der Distanz 1 von Erde und
Sonne.

\begin{center}\fontsize{10}{12}\selectfont\setlength{\tabcolsep}{3pt}
\begin{tabular}{ll|ll|lll|l}
\multicolumn{2}{c|}{1818} & \multicolumn{2}{c|}{AR.} & \multicolumn{3}{c|}{Decl.} & Lichtstärke \\
\hline
April & 6  & $298^\circ$ & $58'$ & $12^\circ$ & $57'$ & N. & 0.396 \\
      & 10 & 297 & 50 & 11 & 2  &    & 0.418 \\
      & 14 & 296 & 27 & 8  & 54 &    & 0.443 \\
      & 18 & 294 & 47 & 6  & 30 &    & 0.471 \\
      & 22 & 292 & 44 & 3  & 48 &    & 0.503 \\
      & 26 & 290 & 17 & 0  & 43 &    & 0.535 \\
      & 30 & 287 & 21 & 2  & 46 & S. & 0.568 \\
Mai   & 4  & 283 & 53 & 6  & 41 &    & 0.596 \\
      & 8  & 279 & 47 & 10 & 58 &    & 0.616 \\
      & 12 & 275 & 4  & 15 & 33 &    & 0.621 \\
\end{tabular}
\end{center}

\par\indent Länger werden wegen des Mondscheins die Beobachtungen wohl nicht fortgesetzt
werden können, und wenn dieser aufhört, wird der Comet für uns wenigstens zu
weit nach Süden fortgerückt sein.

\par\indent Ich selbst habe den Cometen bis jetzt erst einmal beobachtet, und mit einem
Stern 9\textsuperscript{ter} Grösse verglichen, der in der Hist. cél. nicht vorkommt. Ich
behalte mir vor, dessen Ort künftig zu bestimmen; ungefähr wird er
$300^\circ\,3'$ AR., $15^\circ\,9'$ Declination haben.

\begin{center}\fontsize{9}{11}\selectfont\setlength{\tabcolsep}{3pt}\begin{tabular}{lrrrll}
April 1. & $15^{\mathrm u}$ & $45^{\mathrm m}$ & $28^{\mathrm s}$ & Sternzeit, & \shortstack[l]{der Comet folgt dem Stern\\$14''98$ in Zeit}\\
& $15$ & $44$ & $56$ & — & \shortstack[l]{der Comet ist $4'\,22''$\\südlicher als der Stern.}\\
\end{tabular}\end{center}

\par\indent Die Declination möchte wohl beträchtlich weniger genau sein, als die
Rectascension.

\Needspace{8\baselineskip}\par\indent [Handschrift:] Beobachtungen dieses Sterns
\begin{center}\fontsize{10}{12}\selectfont\begin{tabular}{lrl}
1819 August & 27 & $20^{\mathrm u}\;0^{\mathrm m}\;11^{\mathrm s}49$\\
Sept. & 4 & $11.48$\\
& 7 & $11.38$\\
& 24 & $11.30$\\
\end{tabular}\end{center}

\clearpage\fancyhead[L]{\small COMET.}\hypertarget{native-408-part-lower}{}\label{native-408-part-lower}
\begin{center}\large COMET.\end{center}
\begin{center}\hrule\vspace{4pt}\fontsize{10}{12}\selectfont Göttingische gelehrte Anzeigen. Stück 60. Seite 593..594. 1818. April 13.\par\vspace{4pt}\hrule\vspace{4pt}\end{center}
\par\indent Zur Erleichterung der Beobachtung des vom Hrn. \textsc{Pons} Ende vorigen Jahrs
entdeckten Cometen hat Hr. Hofr. \textsc{Gauss} für die noch übrige Zeit
seiner Sichtbarkeit den Lauf desselben im Voraus berechnet; wir glauben durch
die Bekanntmachung dieser Resultate den Beobachtern einen Dienst zu
erweisen. Die parabolischen Elemente, nach denen diese Ephemeride berechnet
ist, gründen sich auf die Beobachtungen des Hrn. Doctor \textsc{Olbers} vom
3. bis 28. März und sind folgende:

\begin{center}\fontsize{10}{12}\selectfont
Zeit des Durchgangs durch die Sonnennähe
\par 1818 Febr. 26. $21^{\mathrm u}\;34^{\mathrm m}\;42^{\mathrm s}$ Zeit im Mer. von Göttingen
\par\smallskip\begin{tabular}{lrr}
Kleinster Abstand von der Sonne & $\cdot\;\cdot\;\cdot\;\cdot$ & $1.20053$\\
Aufsteigender Knoten & $\cdot\;\cdot\;\cdot\;\cdot\;\cdot\;\cdot$ & $70^\circ\;5'\;12''$\\
Neigung der Bahn & $\cdot\;\cdot\;\cdot\;\cdot\;\cdot\;\cdot\;\cdot$ & $90\;0\;0$\\
Länge der Sonnennähe & $\cdot\;\cdot\;\cdot\;\cdot\;\cdot\;\cdot$ & $183\;22\;58$\\
\end{tabular}\end{center}

\SourceLeaf{409}\par\indent Die künftige Verbesserung dieser Elemente nach mehrern Beobachtungen wird erst lehren, ob
die Neigung merklich vom rechten Winkel abweicht, und in welchem Sinn; bis jetzt bleibt
unbestimmt, ob die Bewegung in Beziehung auf die Ekliptik rechtläufig oder rückläufig ist.
Die Zeit der Ephemeride ist $14^{\mathrm u}\,20^{\mathrm m}$ im Meridian von Göttingen. [Die Dphemeride ist
in dem vorhergehenden Aufsatze abgedruckt.]

\par\indent Obgleich der Comet schon weit über seine Sonnennähe hinaus ist, so wird doch sein Licht
wegen der noch abnehmenden Entfernung von der Erde noch immer zunehmen, so dass es am
12. Mai doppelt so gross sein wird als am 13. März; allein da dasselbe an sich so
ausserordentlich schwach ist, so wird der Comet doch nur durch gute Fernröhre zu beobachten
sein.

\par\indent Die auf der hiesigen Sternwarte angestellten Beobachtungen dieses Cometen werden bei einer
andern Gelegenheit bekannt gemacht werden.
